\section{结论}
在 2-bit block-scaled VQ 中，已有组尺度和当前码字邻域 assignment 并非等价适应接口：前者密集、高任务收益且语言漂移较大，后者稀疏、收益较小且漂移较低。把二者置于同一软松弛并未组合优势；naive joint 被尺度端点严格支配，零切换 joint 尺度的退化进一步定位了 soft mixture 补偿与 hard deployment 的失配。证据导出“连续与离散坐标只通过真实 hard checkpoint 交接”的原则，但分阶段修复仍需下一轮实验验证。当前贡献是可复现的现象和机制边界，而不是尚未成立的联合 SOTA。
